\documentclass{beamer}
\usepackage[utf8]{inputenc}
\usepackage[T1]{fontenc}
\usepackage[ngerman]{babel}
\usepackage{amsmath,amssymb,amsthm,mathtools}
\usepackage{bbm}
\usepackage{tikz}
\usetikzlibrary{positioning}
\usepackage{graphicx}
\usepackage{csquotes}
\setbeamertemplate{footline}[frame number]
\setbeamertemplate{navigation symbols}{}
\newcommand{\N}{\mathbbm{N}} %natürliche Zahlen
\newcommand{\Z}{\mathbbm{Z}} %ganze Zahlem
\newcommand{\R}{\mathbb{R}} %relle Zahlen
\newtheorem{satz}{Satz}
\newtheorem{algorithmus}{Algorithmus}
\newtheorem{korollar}{Korollar}
\title{Einfache Graphenalgorithmen}
\subtitle{Proseminar Kombinatorik}
\author{Lennart Meier}
\date{26.05.2026}
\begin{document}
\begin{frame}
\titlepage
\end{frame}
\begin{frame}{Gliederung}
\tableofcontents
\end{frame}
\section{Einleitung}
\begin{frame}{Anwendung}
Graphen gehören zu den wichtigsten Strukturen der Mathematik und Informatik.
\vspace{0.4cm}
Anwendungen:
\begin{itemize}
    \item Straßennetze und Bahnnetze
    \item soziale Netzwerke
    \item Kommunikationsnetze
    \item Suchmaschinen
    \item Projektplanung
\end{itemize}
\vspace{0.4cm}
Zentrale Fragen:
\begin{itemize}
    \item Welche Knoten sind erreichbar?
    \item Wie findet man kürzeste Wege?
    \item Ist ein Graph bipartit?
    \item Enthält ein Digraph Kreise?
\end{itemize}
\end{frame}
\section{Grundbegriffe}
\begin{frame}{Definition eines Graphen}
\begin{definition}
Ein Graph ist ein Paar
\[G=(V,E),\]
wobei
\begin{itemize}
    \item \(V\) eine endliche Menge von Knoten ist,
    \item \(E\) eine Menge von Kanten ist.
\end{itemize}
\end{definition}
\vspace{0.4cm}
Bei ungerichteten Graphen besitzt jede Kante die Form
\[\{u,v\}.\]
Bei gerichteten Graphen verwendet man geordnete Paare
\[(u,v).\]
\end{frame}
\begin{frame}{Beispiel eines Graphen}
\begin{center}
\begin{tikzpicture}[scale=1.2]
\node[circle,draw] (A) at (0,0) {A};
\node[circle,draw] (B) at (2,1) {B};
\node[circle,draw] (C) at (2,-1) {C};
\node[circle,draw] (D) at (4,0) {D};
\draw (A)--(B);
\draw (A)--(C);
\draw (B)--(D);
\draw (C)--(D);
\end{tikzpicture}
\end{center}
\vspace{0.3cm}
\[V=\{A,B,C,D\}\]
\[E=\{\{A,B\},\{A,C\},\{B,D\},\{C,D\}\}\]
\end{frame}
\begin{frame}{Wege und Kreise}
\begin{definition}
\begin{itemize}
    \item Ein Weg ist eine Folge von Knoten, bei der aufeinanderfolgende Knoten durch Kanten verbunden sind.  
    \item Ein Kreis ist ein geschlossener Weg.   
    \item Ein Graph heißt zusammenhängend, wenn zwischen je zwei Knoten ein Weg existiert.
\end{itemize}
\end{definition}
\end{frame}
\section{Graphendurchmusterung}
\begin{frame}{Idee der Graphendurchmusterung}
Ziel:
\begin{itemize}
    \item Bestimme alle von einem Startknoten erreichbaren Knoten.
\end{itemize}
\vspace{0.4cm}
Dabei entstehen:
\begin{itemize}
    \item eine Menge erreichbarer Knoten \(R\)
    \item ein Suchbaum \(F\)
\end{itemize}
\vspace{0.4cm}
Grundlage vieler Algorithmen:
\begin{itemize}
    \item DFS
    \item BFS
    \item Komponentenzerlegung
    \item Kürzeste Wege
\end{itemize}
\end{frame}
\begin{frame}{Algorithmus zur Graphendurchmusterung}
\begin{algorithmus}
Eingabe: ein Graph \(G\), ein Knoten \(r \in V(G)\)
Ausgabe: die Menge \(R \subseteq V(G)\) der von \(r\) aus erreichbaren Knoten und eine Menge \(F \subseteq E(G)\), sodass \((R,F)\) ein Baum beziehungsweise eine Arboreszenz mit Wurzel \(r\) ist.
\begin{align*}
&R \leftarrow \{r\},\quad Q \leftarrow \{r\},\quad F \leftarrow \emptyset\\
&\textbf{while } Q \neq \emptyset \textbf{ do}\\
&\qquad \text{Wähle ein } v \in Q\\
&\qquad \textbf{if } \exists e=(v,w)\in \delta_G^+(v)\\
&\qquad\qquad \text{mit } w\in V(G)\setminus R\\
&\qquad \textbf{then } R\leftarrow R\cup\{w\}\\
&\qquad\qquad Q\leftarrow Q\cup\{w\}\\
&\qquad\qquad F\leftarrow F\cup\{e\}\\
&\qquad \textbf{else } Q\leftarrow Q\setminus\{v\}
\end{align*}
\end{algorithmus}
\end{frame}
\begin{frame}{Korrektheit und Laufzeit}
\begin{satz}
Der Algorithmus arbeitet korrekt und besitzt Laufzeit
\[O(n+m),\]
wobei
\[n=|V(G)|,\qquad m=|E(G)|.\]
\end{satz}
\vspace{0.4cm}
\begin{proof}
\begin{itemize}
    \item Jeder Knoten wird höchstens einmal eingefügt.
    \item Jede Kante wird höchstens zweimal betrachtet.
    \item Dadurch ergibt sich lineare Laufzeit.
    \item Am Ende enthält \(R\) genau alle erreichbaren Knoten.
\end{itemize}
\end{proof}
\end{frame}
\section{Tiefensuche}
\begin{frame}{Definition der Tiefensuche}
\begin{definition}
Wählt man in der Menge \(Q\) stets den zuletzt hinzugefügten Knoten aus, so erhält man die Tiefensuche beziehungsweise DFS (\emph{Depth First Search}).
\end{definition}
\vspace{0.5cm}
Prinzip:
\begin{quote}
Gehe möglichst tief in den Graphen.
\end{quote}
\vspace{0.4cm}
DFS verwendet einen Stack.
\end{frame}
\begin{frame}{Beispiel zur Tiefensuche}
\begin{center}
\begin{tikzpicture}[scale=1.1]
\node[circle,draw] (A) at (0,0) {A};
\node[circle,draw] (B) at (2,1) {B};
\node[circle,draw] (C) at (2,-1) {C};
\node[circle,draw] (D) at (4,1) {D};
\node[circle,draw] (E) at (4,-1) {E};
\draw (A)--(B);
\draw (A)--(C);
\draw (B)--(D);
\draw (C)--(E);
\end{tikzpicture}
\end{center}
\vspace{0.3cm}
Mögliche DFS-Reihenfolge:
\[A,B,D,C,E\]
\end{frame}
\begin{frame}{Anwendungen der Tiefensuche}
DFS wird verwendet für:
\begin{itemize}
    \item Zusammenhangsprobleme
    \item Kreissuche
    \item Komponentenzerlegung
    \item topologische Sortierung
    \item Labyrinthsuche
    \item Dateisysteme
\end{itemize}
\end{frame}
\section{Breitensuche}
\begin{frame}{Definition der Breitensuche}
\begin{definition}
Wählt man in der Menge \(Q\) stets den zuerst hinzugefügten Knoten aus, so erhält man die Breitensuche beziehungsweise BFS (\emph{Breadth First Search}).
\end{definition}
\vspace{0.5cm}
Prinzip:
\begin{quote}
Durchsuche den Graphen schichtweise.
\end{quote}
\vspace{0.4cm}
BFS verwendet eine Queue.
\end{frame}
\begin{frame}{Beispiel zur Breitensuche}
\begin{center}
\begin{tikzpicture}[scale=1.1]
\node[circle,draw] (A) at (0,0) {A};
\node[circle,draw] (B) at (2,1.5) {B};
\node[circle,draw] (C) at (2,0) {C};
\node[circle,draw] (D) at (2,-1.5) {D};
\node[circle,draw] (E) at (4,1) {E};
\node[circle,draw] (F) at (4,-1) {F};
\draw (A)--(B);
\draw (A)--(C);
\draw (A)--(D);
\draw (B)--(E);
\draw (D)--(F);
\end{tikzpicture}
\end{center}
\vspace{0.3cm}
BFS-Ebenen:
\[\{A\},\{B,C,D\},\{E,F\}\]
\end{frame}
\begin{frame}{Kürzeste Wege}
Für zwei Knoten \(v,w\) bezeichnet
\[dist(v,w)\]
die Länge eines kürzesten Weges.
\vspace{0.5cm}
\begin{satz}
Jeder BFS-Baum enthält einen kürzesten Weg vom Startknoten \(r\) zu allen erreichbaren Knoten.
\end{satz}
\end{frame}
\begin{frame}{Beweisidee}
\begin{proof}
\begin{itemize}
    \item BFS arbeitet schichtweise.
    \item Zunächst werden alle Knoten mit Distanz \(1\) besucht.
    \item Danach alle Knoten mit Distanz \(2\) usw.
    \item Ein Knoten wird daher erstmals über einen kürzesten Weg erreicht.
\end{itemize}
\end{proof}
\end{frame}
\begin{frame}{Laufzeit der Breitensuche}
\begin{satz}
Die Werte
\[dist(r,v)\]
können für alle Knoten \(v\) in linearer Zeit bestimmt werden.
\end{satz
\vspace{0.4cm}
\begin{proof}
\begin{itemize}
    \item Jeder Knoten wird höchstens einmal besucht.
    \item Jede Kante wird höchstens zweimal betrachtet.
    \item Daher beträgt die Laufzeit
    \[ O(n+m).\]
\end{itemize}
\end{proof}
\end{frame}
\section{Bipartite Graphen}
\begin{frame}{Definition bipartiter Graphen}
\begin{definition}
Eine Bipartition eines Graphen besteht aus zwei disjunkten Mengen
\[A,B\]
mit
\[A\cup B=V(G),\]
sodass jede Kante genau einen Endknoten in \(A\) besitzt.
\end{definition}
\vspace{0.4cm}
Ein Graph heißt bipartit, falls er eine Bipartition besitzt.
\end{frame}
\begin{frame}{Beispiel eines bipartiten Graphen}
\begin{center}
\begin{tikzpicture}[scale=1]
\node[circle,draw] (A) at (0,2) {A};
\node[circle,draw] (B) at (0,0) {B};
\node[circle,draw] (C) at (4,3) {C};
\node[circle,draw] (D) at (4,1.5) {D};
\node[circle,draw] (E) at (4,0) {E};
\draw (A)--(C);
\draw (A)--(D);
\draw (B)--(D);
\draw (B)--(E);
\end{tikzpicture}
\end{center}
\[A=\{A,B\},
\qquad
B=\{C,D,E\}\]
\end{frame}
\begin{frame}{Charakterisierung}
\begin{definition}
Ein ungerader Kreis ist ein Kreis ungerader Länge.
\end{definition}
\vspace{0.4cm}
\begin{satz}[König]
Ein ungerichteter Graph ist genau dann bipartit, wenn er keinen ungeraden Kreis enthält.
\end{satz}
\end{frame}
\begin{frame}{Beweisidee}
\begin{proof}
\begin{itemize}
    \item In bipartiten Graphen wechselt man entlang jeder Kante die Klasse.
    \item Deshalb besitzen alle Kreise gerade Länge.
    \item umgekehrt zerlegt BFS die Knoten in gerade und ungerade Ebenen.
    \item Eine Kante innerhalb einer Ebene würde einen ungeraden Kreis erzeugen.
\end{itemize}
\end{proof}
\end{frame}
\section{Azyklische Digraphen}
\begin{frame}{Definition azyklischer Digraphen}
\begin{definition}
Ein Digraph heißt azyklisch, falls er keinen gerichteten Kreis enthält.
\end{definition}
\vspace{0.5cm}
Solche Graphen heißen auch:
\[\text{DAG}\]
(\emph{Directed Acyclic Graphs})
\end{frame}
\begin{frame}{Topologische Ordnung}
\begin{definition}
Eine topologische Ordnung eines Digraphen ist eine Reihenfolge
\[v_1,\dots,v_n,\]
sodass für jede Kante
\[(v_i,v_j)\]
gilt:
\[i<j.\]
\end{definition}
\end{frame}
\begin{frame}{Lemma}
\begin{lemma}
Ist \(G\) ein Digraph mit
\[\delta^+(v)\neq\emptyset\]
für alle \(v\in V(G)\), so kann man in
\[O(|V(G)|)\]
Zeit einen Kreis finden.
\end{lemma}
\end{frame}
\begin{frame}{Beweisidee}
\begin{proof}
\begin{itemize}
    \item Starte in einem beliebigen Knoten.
    \item Folge immer einer ausgehenden Kante
    \item Da der Graph endlich ist, wird irgendwann ein Knoten doppelt besucht.
    \item Dadurch entsteht ein Kreis.
\end{itemize}
\end{proof}
\end{frame}
\begin{frame}{Topologische Ordnung und DAG}
\begin{satz}
Ein Digraph besitzt genau dann eine topologische Ordnung, wenn er azyklisch ist.
\end{satz}
\end{frame}
\begin{frame}{Beweisidee}
\begin{proof}
\begin{itemize}
    \item In einer topologischen Ordnung wächst die Nummerierung entlang jeder Kante.
    \item Deshalb kann kein gerichteter Kreis existieren.
    \item Umgekehrt entfernt man wiederholt Knoten ohne ausgehende Kanten.
    \item Dadurch entsteht eine topologische Ordnung.
\end{itemize}
\end{proof}
\end{frame}
\section{Fazit}
\begin{frame}{Zusammenfassung}
\begin{itemize}
    \item Graphendurchmusterung ist Grundlage vieler Verfahren.
    \item DFS und BFS besitzen Laufzeit
    \[O(n+m).\]
    \item BFS berechnet kürzeste Wege.
    \item Bipartite Graphen enthalten keine ungeraden Kreise.
    \item DAGs besitzen topologische Ordnungen.
\end{itemize}
\end{frame}
\begin{frame}{Vielen Dank!}
Vielen Dank für eure Aufmerksamkeit!
\vspace{1cm}
Fragen?
\end{frame}
\end{document}
